\phantomsection\section*{ЗАКЛЮЧЕНИЕ}\addcontentsline{toc}{section}{ЗАКЛЮЧЕНИЕ}

В результате курсовой работы разработано программное обеспечение для визуализации сцены из геометрических объектов и неравномерного тумана методом трассировки путей.

В ходе работы были решены следующие задачи:

\begin{itemize}
    \item Проанализированы существующие методы моделирования тумана, проведён сравнительный анализ алгоритмов рендеринга. Для реализации выбраны экспоненциальная модель тумана и метод трассировки путей.
    \item Разработан алгоритм генерации изображения с поддержкой взаимодействия лучей с туманом, включающий стохастическое определение расстояния до взаимодействия, вычисление неоднородной плотности с использованием шума Перлина и экспоненциального спада по высоте.
    \item Реализован алгоритм генерации изображения методом трассировки путей.
    \item Реализована модель объёмного тумана с переменной плотностью на основе шума Перлина, создающая реалистичные турбулентные структуры и концентрацию у поверхности земли.
    \item Проведено исследование зависимости производительности от глубины рекурсии и количества лучей на пиксель. Установлена линейная зависимость от обоих параметров. Определены оптимальные параметры для интерактивной работы и финального рендеринга.
\end{itemize}

Разработанное программное обеспечение обеспечивает физически корректную визуализацию объёмного тумана с глобальным освещением и поддерживает редактирование параметров через графический интерфейс.
Таким образом, все поставленные задачи решены, цель работы достигнута.

\clearpage
